\documentclass[12pt,a4paper]{article}

% ---- Encoding & Font ----
\usepackage{iftex}
\ifPDFTeX
  \usepackage[utf8]{inputenc}
  \usepackage[T1]{fontenc}
\else
  \usepackage{fontspec}  % XeTeX/LuaTeX (tectonic): Unicode nativo (·, —, ¿, ª, §)
\fi
\usepackage[spanish]{babel}
\usepackage{csquotes}

% ---- Page layout ----
\usepackage[top=2.5cm, bottom=2.5cm, left=2.5cm, right=2.5cm]{geometry}
\usepackage{setspace}
\onehalfspacing

% ---- Math ----
\usepackage{amsmath, amssymb, amsthm}
\usepackage{mathtools}
\usepackage{physics}
\usepackage{esint}  % \oiint

% ---- Graphics ----
\usepackage{graphicx}
\usepackage{tikz}
\usepackage{float}
\usepackage{caption}

% ---- Tables ----
\usepackage{array}
\usepackage{booktabs}
\usepackage{tabularx}
\usepackage{colortbl}

% ---- Colours ----
\usepackage{xcolor}
\definecolor{notebg}{HTML}{FFF3CD}
\definecolor{noteborder}{HTML}{FFC107}
\definecolor{boxred}{HTML}{F8D7DA}
\definecolor{boxredborder}{HTML}{F5C6CB}

% ---- Boxes ----
\usepackage{mdframed}
\usepackage{tcolorbox}
\tcbuselibrary{most}

\newtcolorbox{keybox}[1][]{
  colback=blue!5,
  colframe=blue!70!black,
  arc=4pt,
  boxrule=1pt,
  fonttitle=\bfseries,
  title={#1}
}

\newtcolorbox{warnbox}[1][]{
  colback=boxred,
  colframe=boxredborder,
  coltitle=black!75,
  arc=4pt,
  boxrule=1pt,
  fonttitle=\bfseries,
  title={#1}
}

\newtcolorbox{notebox}[1][]{
  colback=notebg,
  colframe=noteborder,
  coltitle=black!75,
  arc=4pt,
  boxrule=1pt,
  fonttitle=\bfseries,
  title={#1}
}

% ---- Diagramas (gráficas y esquemas) ----
\usepackage{pgfplots}
\pgfplotsset{compat=1.18}
\usetikzlibrary{babel}  % babel español activa `>`; esto evita romper las flechas `->`
% courses/ElecMag2024/Clases/assets/tikzstyles.tex
% ---------------------------------------------------------------------------
% Estilos compartidos para los diagramas del curso Electromagnetismo (2024).
% Fuente única del "look" visual (colores, grosores, escalas) para que todas
% las figuras (TikZ / pgfplots / CircuiTikz) sean consistentes entre clases.
%
% Es una COPIA de courses/Fisica3-2015/Clases/assets/tikzstyles.tex: mismo
% dominio, misma paleta. Copia y no symlink para que cada curso sea
% autocontenido y la paleta de Física III pueda evolucionar aparte.
%
% Uso: la clase debe cargar antes `tikz` y `pgfplots` (y `circuitikz` si la
%      clase tiene circuitos), y luego % courses/ElecMag2024/Clases/assets/tikzstyles.tex
% ---------------------------------------------------------------------------
% Estilos compartidos para los diagramas del curso Electromagnetismo (2024).
% Fuente única del "look" visual (colores, grosores, escalas) para que todas
% las figuras (TikZ / pgfplots / CircuiTikz) sean consistentes entre clases.
%
% Es una COPIA de courses/Fisica3-2015/Clases/assets/tikzstyles.tex: mismo
% dominio, misma paleta. Copia y no symlink para que cada curso sea
% autocontenido y la paleta de Física III pueda evolucionar aparte.
%
% Uso: la clase debe cargar antes `tikz` y `pgfplots` (y `circuitikz` si la
%      clase tiene circuitos), y luego % courses/ElecMag2024/Clases/assets/tikzstyles.tex
% ---------------------------------------------------------------------------
% Estilos compartidos para los diagramas del curso Electromagnetismo (2024).
% Fuente única del "look" visual (colores, grosores, escalas) para que todas
% las figuras (TikZ / pgfplots / CircuiTikz) sean consistentes entre clases.
%
% Es una COPIA de courses/Fisica3-2015/Clases/assets/tikzstyles.tex: mismo
% dominio, misma paleta. Copia y no symlink para que cada curso sea
% autocontenido y la paleta de Física III pueda evolucionar aparte.
%
% Uso: la clase debe cargar antes `tikz` y `pgfplots` (y `circuitikz` si la
%      clase tiene circuitos), y luego \input{../assets/tikzstyles.tex}
% (compilar desde el directorio de la clase para que la ruta relativa resuelva).
% ---------------------------------------------------------------------------

% ---- Paleta (alineada con las cajas del preámbulo: keybox/notebox/warnbox) ----
\definecolor{figblue}{HTML}{1F4E79}   % azul principal (líneas, keybox)
\definecolor{figred}{HTML}{C0392B}    % rojo de énfasis (warnbox)
\definecolor{figamber}{HTML}{B8860B}  % ámbar (notebox)
\definecolor{figgray}{HTML}{555555}   % auxiliares, ejes, guías

% ---- CircuiTikz: escala y grosor de los componentes ----
% Guarda \ifdefined: la electrostática (Clases 1-14) no carga `circuitikz`, y
% sin la guarda el \input revienta con `Undefined control sequence \ctikzset`.
% Cuando el curso llegue a circuitos, el bloque se activa solo.
\ifdefined\ctikzset
  \ctikzset{
    bipoles/thickness=1,
    resistors/scale=0.9,
    capacitors/scale=0.9,
  }
\fi

% ---- TikZ: estilos reutilizables ----
% Nota: las puntas se declaran como `-{Latex[...]}` y no como `->`. La forma
% con llaves NO contiene un `>` literal, así que es inmune al `>` activo que
% introduce babel español (ver CLAUDE.md §7.3.3).
\usetikzlibrary{arrows.meta,calc,angles,quotes}

\tikzset{
  figline/.style={figblue, line width=0.9pt},
  figaux/.style={figgray, dashed, line width=0.6pt},
  % vectores
  figvec/.style={figblue, line width=0.9pt, -{Latex[length=2.2mm,width=1.6mm]}},
  figvecres/.style={figred, line width=1.3pt, -{Latex[length=2.6mm,width=1.9mm]}},
  figvecaux/.style={figgray, line width=0.7pt, -{Latex[length=1.8mm,width=1.3mm]}},
  figaxis/.style={figgray, line width=0.6pt, -{Latex[length=2mm,width=1.4mm]}},
  % cargas puntuales
  figcarga/.style={circle, inner sep=0pt, minimum size=5.4pt, draw=none},
  figpos/.style={figcarga, fill=figred},
  figneg/.style={figcarga, fill=figblue},
  figneutra/.style={figcarga, fill=figgray},
  % etiquetas
  figlbl/.style={font=\small},
  figlblaux/.style={font=\footnotesize, figgray},
}

% ---- pgfplots: estilo base de las gráficas ----
\pgfplotsset{
  emfig/.style={
    width=11cm, height=6.5cm,
    axis lines=left,
    xlabel style={font=\small},
    ylabel style={font=\small},
    tick label style={font=\footnotesize},
    every axis plot/.append style={line width=1pt},
    grid=major,
    grid style={figgray!20},
    legend cell align=left,
    legend style={font=\footnotesize, draw=figgray!40},
  },
}

% (compilar desde el directorio de la clase para que la ruta relativa resuelva).
% ---------------------------------------------------------------------------

% ---- Paleta (alineada con las cajas del preámbulo: keybox/notebox/warnbox) ----
\definecolor{figblue}{HTML}{1F4E79}   % azul principal (líneas, keybox)
\definecolor{figred}{HTML}{C0392B}    % rojo de énfasis (warnbox)
\definecolor{figamber}{HTML}{B8860B}  % ámbar (notebox)
\definecolor{figgray}{HTML}{555555}   % auxiliares, ejes, guías

% ---- CircuiTikz: escala y grosor de los componentes ----
% Guarda \ifdefined: la electrostática (Clases 1-14) no carga `circuitikz`, y
% sin la guarda el \input revienta con `Undefined control sequence \ctikzset`.
% Cuando el curso llegue a circuitos, el bloque se activa solo.
\ifdefined\ctikzset
  \ctikzset{
    bipoles/thickness=1,
    resistors/scale=0.9,
    capacitors/scale=0.9,
  }
\fi

% ---- TikZ: estilos reutilizables ----
% Nota: las puntas se declaran como `-{Latex[...]}` y no como `->`. La forma
% con llaves NO contiene un `>` literal, así que es inmune al `>` activo que
% introduce babel español (ver CLAUDE.md §7.3.3).
\usetikzlibrary{arrows.meta,calc,angles,quotes}

\tikzset{
  figline/.style={figblue, line width=0.9pt},
  figaux/.style={figgray, dashed, line width=0.6pt},
  % vectores
  figvec/.style={figblue, line width=0.9pt, -{Latex[length=2.2mm,width=1.6mm]}},
  figvecres/.style={figred, line width=1.3pt, -{Latex[length=2.6mm,width=1.9mm]}},
  figvecaux/.style={figgray, line width=0.7pt, -{Latex[length=1.8mm,width=1.3mm]}},
  figaxis/.style={figgray, line width=0.6pt, -{Latex[length=2mm,width=1.4mm]}},
  % cargas puntuales
  figcarga/.style={circle, inner sep=0pt, minimum size=5.4pt, draw=none},
  figpos/.style={figcarga, fill=figred},
  figneg/.style={figcarga, fill=figblue},
  figneutra/.style={figcarga, fill=figgray},
  % etiquetas
  figlbl/.style={font=\small},
  figlblaux/.style={font=\footnotesize, figgray},
}

% ---- pgfplots: estilo base de las gráficas ----
\pgfplotsset{
  emfig/.style={
    width=11cm, height=6.5cm,
    axis lines=left,
    xlabel style={font=\small},
    ylabel style={font=\small},
    tick label style={font=\footnotesize},
    every axis plot/.append style={line width=1pt},
    grid=major,
    grid style={figgray!20},
    legend cell align=left,
    legend style={font=\footnotesize, draw=figgray!40},
  },
}

% (compilar desde el directorio de la clase para que la ruta relativa resuelva).
% ---------------------------------------------------------------------------

% ---- Paleta (alineada con las cajas del preámbulo: keybox/notebox/warnbox) ----
\definecolor{figblue}{HTML}{1F4E79}   % azul principal (líneas, keybox)
\definecolor{figred}{HTML}{C0392B}    % rojo de énfasis (warnbox)
\definecolor{figamber}{HTML}{B8860B}  % ámbar (notebox)
\definecolor{figgray}{HTML}{555555}   % auxiliares, ejes, guías

% ---- CircuiTikz: escala y grosor de los componentes ----
% Guarda \ifdefined: la electrostática (Clases 1-14) no carga `circuitikz`, y
% sin la guarda el \input revienta con `Undefined control sequence \ctikzset`.
% Cuando el curso llegue a circuitos, el bloque se activa solo.
\ifdefined\ctikzset
  \ctikzset{
    bipoles/thickness=1,
    resistors/scale=0.9,
    capacitors/scale=0.9,
  }
\fi

% ---- TikZ: estilos reutilizables ----
% Nota: las puntas se declaran como `-{Latex[...]}` y no como `->`. La forma
% con llaves NO contiene un `>` literal, así que es inmune al `>` activo que
% introduce babel español (ver CLAUDE.md §7.3.3).
\usetikzlibrary{arrows.meta,calc,angles,quotes}

\tikzset{
  figline/.style={figblue, line width=0.9pt},
  figaux/.style={figgray, dashed, line width=0.6pt},
  % vectores
  figvec/.style={figblue, line width=0.9pt, -{Latex[length=2.2mm,width=1.6mm]}},
  figvecres/.style={figred, line width=1.3pt, -{Latex[length=2.6mm,width=1.9mm]}},
  figvecaux/.style={figgray, line width=0.7pt, -{Latex[length=1.8mm,width=1.3mm]}},
  figaxis/.style={figgray, line width=0.6pt, -{Latex[length=2mm,width=1.4mm]}},
  % cargas puntuales
  figcarga/.style={circle, inner sep=0pt, minimum size=5.4pt, draw=none},
  figpos/.style={figcarga, fill=figred},
  figneg/.style={figcarga, fill=figblue},
  figneutra/.style={figcarga, fill=figgray},
  % etiquetas
  figlbl/.style={font=\small},
  figlblaux/.style={font=\footnotesize, figgray},
}

% ---- pgfplots: estilo base de las gráficas ----
\pgfplotsset{
  emfig/.style={
    width=11cm, height=6.5cm,
    axis lines=left,
    xlabel style={font=\small},
    ylabel style={font=\small},
    tick label style={font=\footnotesize},
    every axis plot/.append style={line width=1pt},
    grid=major,
    grid style={figgray!20},
    legend cell align=left,
    legend style={font=\footnotesize, draw=figgray!40},
  },
}


% ---- Hyperlinks & TOC ----
\usepackage[hidelinks]{hyperref}
\usepackage{bookmark}

% ---- Enumerate ----
\usepackage{enumitem}

% ---- Miscellaneous ----
\usepackage{icomma}
\usepackage{siunitx}
\usepackage{textcomp}
\usepackage[bottom]{footmisc}

% ---- Title ----
\title{\textbf{Clase 4: Conductores en equilibrio y desarrollo multipolar} \\[0.3em]
        \large{Del dipolo eléctrico al desarrollo de Taylor de $1/|\vec r-\vec r\,'|$}}
\author{Transcripto y expandido --- Curso Electromagnetismo (OpenFING)}
\date{}

\begin{document}

\maketitle
\thispagestyle{empty}
\newpage

\tableofcontents
\newpage

\section{Repaso: de Gauss a Coulomb}

La clase anterior había mostrado que la ley de Coulomb \textbf{implica} la ley de
Gauss. La clase abre cerrando el círculo: la implicación también vale al revés, y
el caso de la carga puntual lo exhibe.

Sea una carga puntual $Q>0$ en el origen. El argumento no arranca por la cuenta
sino por la \textbf{simetría}: si la carga se piensa puntual, todas las
direcciones son equivalentes, de modo que

\begin{itemize}[nosep]
  \item $\vec E$ es \textbf{radial} en todo punto, y
  \item $|\vec E|$ es el mismo en dos puntos a igual distancia de la carga.
\end{itemize}

\begin{notebox}[El orden importa]
Dirección y sentido se determinan \textbf{a priori}, antes de aplicar Gauss. La
ley de Gauss sólo aporta el módulo. Este orden ---simetría primero, ley
después--- es el que hace útil a Gauss para calcular.
\end{notebox}

Con una superficie gaussiana esférica de radio $r$ centrada en la carga, el flujo
saliente sale de la integral porque el módulo es constante sobre ella:

\[
\oiint_S \vec E \cdot d\vec S \;=\; |\vec E| \oiint_S dS \;=\; 4\pi r^2\,|\vec E|
\;=\; \frac{Q}{\varepsilon_0}
\]

\[
\boxed{\;|\vec E| = \frac{1}{4\pi\varepsilon_0}\frac{Q}{r^2}\;}
\]

que es Coulomb, con el sentido saliente si $Q>0$ y entrante si $Q<0$.

\section{Conductores en equilibrio}

\subsection{Campo y densidad de carga en el interior}

Sea un \textbf{conductor en equilibrio} (un metal, aunque el argumento no usa que
sea metálico).

\textbf{Primero, $\vec E = 0$ adentro.} El argumento es por el absurdo y es
físico, no matemático: si el campo interior no fuera nulo, las cargas se
moverían. Si hay carga neta, se movería la carga neta; si no la hay, el material
se polarizaría ---las positivas hacia un lado, las negativas hacia el otro---.
Cualquiera de las dos cosas contradice estar en equilibrio.

\textbf{Segundo, $\rho = 0$ adentro.} Esto ya no es un argumento físico
independiente: es una consecuencia inmediata de la forma diferencial de Gauss,
que es justamente la herramienta nueva del curso. De

\[
\div \vec E = \frac{\rho}{\varepsilon_0}
\qquad\Longrightarrow\qquad
\rho = \varepsilon_0\,\div\vec E = 0
\]

porque la divergencia de un campo idénticamente nulo es nula.

\begin{keybox}[Conductor en equilibrio]
$\vec E = 0$ y $\rho = 0$ en el interior. La carga neta, si la hay, sólo puede
estar en \textbf{alguna} superficie.
\end{keybox}

\begin{notebox}[Contraste de método con Física III]
Allá esto se argumenta con superficies gaussianas encogidas; acá sale en un
renglón de la forma diferencial. Es exactamente la ganancia que el curso viene a
capitalizar.
\end{notebox}

\subsection{Conductor con hueco: la carga vive en el borde exterior}

Considérese ahora un conductor en equilibrio con un \textbf{hueco interior sin
cargas} (puede haber carga en el conductor, pero no dentro del hueco). Lo
anterior dice que la carga está en algún borde, pero ahora hay \textbf{dos}
bordes: el exterior y el que rodea el hueco. La propiedad extra es que
\textbf{está toda en el exterior}.

\begin{figure}[H]
  \centering
  % Clase 4 — conductor en equilibrio con hueco sin cargas.
% (a) gaussiana S pegada al hueco, del lado del material: E=0 sobre S => Q_int=0.
% (b) por qué la carga neta nula no basta: parches +/- compensados y la curva C
%     cerrada que los descarta (circulacion no nula).
\begin{tikzpicture}[scale=1]

% ---------------- panel (a) ----------------
\begin{scope}
  \fill[figgray!18] (0,0) circle (2.05);
  \draw[figline] (0,0) circle (2.05);
  \fill[white] plot[smooth cycle, tension=0.9]
    coordinates {(-0.75,0.55) (0.15,0.9) (0.9,0.3) (0.6,-0.65) (-0.4,-0.6)};
  \draw[figline] plot[smooth cycle, tension=0.9]
    coordinates {(-0.75,0.55) (0.15,0.9) (0.9,0.3) (0.6,-0.65) (-0.4,-0.6)};
  % superficie gaussiana S, pegada al hueco pero del lado del material
  \draw[figred, dashed, line width=0.9pt] plot[smooth cycle, tension=0.9]
    coordinates {(-1.0,0.75) (0.2,1.18) (1.15,0.4) (0.82,-0.9) (-0.6,-0.85)};
  \node[figlbl, figred, fill=white, inner sep=1.5pt] at (1.02,1.28) {$S$};
  \node[figlbl] at (0.08,0.1) {hueco};
  \node[figlblaux, fill=figgray!18, inner sep=1.5pt] at (0,-1.5) {$\vec E=0,\ \rho=0$};
  \node[figlbl] at (0,-2.55) {conductor en equilibrio};
  \node[figlbl] at (0,-3.15) {(a) $\vec E=0$ sobre $S\ \Rightarrow\ Q_{\text{sup. int}}=0$};
\end{scope}

% ---------------- panel (b) ----------------
\begin{scope}[xshift=7.4cm]
  \fill[figgray!18] (0,0) circle (2.05);
  \draw[figline] (0,0) circle (2.05);
  \fill[white] plot[smooth cycle, tension=0.85]
    coordinates {(-1.3,0.5) (-0.1,0.85) (1.2,0.5) (1.25,-0.2) (-0.1,-0.55) (-1.3,-0.2)};
  \draw[figline] plot[smooth cycle, tension=0.85]
    coordinates {(-1.3,0.5) (-0.1,0.85) (1.2,0.5) (1.25,-0.2) (-0.1,-0.55) (-1.3,-0.2)};
  % parches de carga sobre la pared del hueco
  \foreach \yy in {-0.1,0.06,0.22}{
    \node[figpos] at (-1.16,\yy) {};
    \node[figneg] at (1.11,\yy) {};
  }
  \node[figlbl, figred, fill=figgray!18, inner sep=1.5pt] at (-1.38,1.08) {$\sigma>0$};
  \node[figlbl, figblue, fill=figgray!18, inner sep=1.5pt] at (1.38,1.08) {$\sigma<0$};
  \draw[figvecaux] (-1.3,0.88) -- (-1.18,0.4);
  \draw[figvecaux] (1.3,0.88) -- (1.18,0.4);
  % campo dentro del hueco: de + a -
  \draw[figvec] (-0.88,0.08) -- (0.78,0.08);
  \node[figlbl] at (-0.05,0.36) {$\vec E$};
  % curva cerrada C: tramo por el hueco y retorno por el material
  \draw[figred, line width=1pt, dashed]
    (-0.95,-0.08) -- (0.82,-0.08) -- (0.82,-1.5) -- (-0.95,-1.5) -- cycle;
  \node[figlbl, figred, fill=figgray!18, inner sep=1.5pt] at (-0.05,-1.5) {$C$};
  \node[figlblaux, fill=figgray!18, inner sep=1.5pt] at (-0.05,-0.9) {aqu\'i $\vec E=0$};
  \node[figlbl] at (0,-2.55) {carga neta nula, pero $\sigma\neq0$};
  \node[figlbl] at (0,-3.15) {(b) $\oint_C\vec E\cdot d\vec\ell\neq0$: imposible};
\end{scope}

\end{tikzpicture}

  \caption*{\small\itshape La gaussiana $S$ del panel (a) prueba que la carga neta
  interna es nula; el panel (b) muestra por qué eso no alcanza y qué configuración
  hay que descartar con la circulación.}
\end{figure}

\textbf{Paso 1 --- la carga neta en la superficie interna es cero.} Se toma una
superficie gaussiana $S$ pegada al hueco pero \textbf{del lado del material}, de
modo que encierre a la superficie interna. Como $S$ está enteramente dentro del
conductor, $\vec E = 0$ sobre ella, y por lo tanto su flujo es nulo. Gauss da
entonces

\[
Q_{\text{encerrada}} = 0
\qquad\Longrightarrow\qquad
Q_{\text{sup. interna}} = 0
\]

(usando que dentro del hueco no hay carga).

\textbf{Paso 2 --- y no sólo la neta: la densidad es cero punto a punto.} Acá
está la sutileza que el docente insiste en no saltear. El paso 1 no excluye una
configuración con \textbf{zonas de $\sigma>0$ y zonas de $\sigma<0$ sobre la
superficie interna que se compensen}. Para descartarla no alcanza Gauss: hay que
usar la otra propiedad del caso electrostático, que la circulación del campo es
nula sobre \textbf{cualquier} curva cerrada.

Si existiera esa configuración de parches, habría un campo eléctrico
\textbf{dentro del hueco}, yendo de las zonas positivas a las negativas. Se toma
entonces una curva cerrada que pase por el hueco (donde $\vec E \neq 0$ y va en el
sentido del recorrido) y se cierre por dentro del material (donde $\vec E = 0$).
El tramo por el material no aporta y el tramo por el hueco aporta positivo, de
modo que

\[
\oint_C \vec E\cdot d\vec\ell > 0
\]

lo que contradice la electrostática. Luego la configuración de parches es
imposible.

\begin{warnbox}[El argumento intuitivo no alcanza]
Decir «las cargas positivas y negativas se atraen, así que se neutralizan»
\textbf{no es una demostración}. El docente construye a propósito un
contraejemplo aparente: un hueco de forma alargada con carga positiva en un
extremo y negativa en el otro parece localmente estable ---«los electrones no van
a querer dar la vuelta»---. Sí dan la vuelta, pero eso no es obvio desde el
argumento de atracción. Lo que cierra el caso es la circulación nula, que es
global.
\end{warnbox}

\subsection{Campo justo afuera: componente tangencial nula}

Ahora el campo \textbf{inmediatamente afuera} del conductor. El resultado
conocido es que es normal a la superficie; la pregunta es cómo se prueba. Se
prueba en dos pasos, y el primero \textbf{no} usa Gauss.

Se toma una \textbf{curva cerrada pequeña} $C$, rectangular y achatada, con un
lado justo afuera de la superficie y el lado opuesto justo adentro del material
(la superficie se supone suave; si tiene aristas el argumento no vale). Como la
curva es chica y achatada, los lados cortos no contribuyen, y de los dos lados
largos:

\begin{itemize}[nosep]
  \item el interior no aporta, porque adentro $\vec E = 0$;
  \item el exterior aporta sólo su \textbf{componente tangencial}, porque el
        producto escalar $\vec E\cdot d\vec\ell$ proyecta sobre la dirección del
        recorrido.
\end{itemize}

Entonces $\oint_C \vec E\cdot d\vec\ell = E_t\,\Delta\ell = 0$, y como
$\Delta\ell \neq 0$,

\[
\boxed{\;E_t = 0 \quad\text{justo afuera del conductor}\;}
\]

Como la dirección tangencial elegida era arbitraria, \textbf{toda} componente
tangencial se anula.

\begin{notebox}[Un error del Resnick]
El docente señala que esta demostración está mal hecha en el Resnick ---uno de
los pocos errores que dice haberle encontrado al libro---.
\end{notebox}

\subsection{Campo justo afuera: componente normal}

Descartada la componente tangencial, sólo queda la normal, y ahora sí entra
Gauss. Se toma una superficie gaussiana con forma de \textbf{cilindro chato y
chico} («cilindrito») a caballo del borde: una tapa afuera, la otra adentro del
material, y altura despreciable.

\begin{figure}[H]
  \centering
  % Clase 4 — campo justo afuera de un conductor, en dos pasos.
% (a) curva cerrada C achatada a caballo del borde: circulacion nula => E_t=0.
% (b) cilindrito gaussiano chato: solo aporta la tapa exterior => E_n=sigma/eps0.
\begin{tikzpicture}[scale=1]

% ---------------- panel (a) ----------------
\begin{scope}
  \fill[figgray!18] (-2.1,-1.5) rectangle (2.1,0);
  \draw[figline] (-2.1,0) -- (2.1,0);
  \node[figlblaux] at (0,-1.15) {conductor: $\vec E=0$};
  % curva C achatada, a caballo del borde
  \draw[figred, line width=1pt] (-0.9,0.3) -- (0.9,0.3) -- (0.9,-0.3) -- (-0.9,-0.3) -- cycle;
  \node[figlbl, figred] at (1.2,0.3) {$C$};
  % componente tangencial exterior, sobre el lado de afuera
  \draw[figvec] (-0.45,0.3) -- (0.5,0.3);
  \node[figlbl, fill=white, inner sep=1.5pt] at (0.02,0.56) {$E_t$};
  % cota del lado largo
  \draw[figaux] (-0.9,0.86) -- (-0.9,1.24);
  \draw[figaux] (0.9,0.86) -- (0.9,1.24);
  \draw[figvecaux] (-0.9,1.06) -- (0.9,1.06);
  \draw[figvecaux] (0.9,1.06) -- (-0.9,1.06);
  \node[figlbl, fill=white, inner sep=1.5pt] at (0,1.06) {$\Delta\ell$};
  \node[figlbl] at (0,-2.25) {(a) $\oint_C\vec E\cdot d\vec\ell=E_t\,\Delta\ell=0$};
  \node[figlbl] at (0,-2.85) {$\Rightarrow\ E_t=0$ (no usa Gauss)};
\end{scope}

% ---------------- panel (b) ----------------
\begin{scope}[xshift=7.4cm]
  \fill[figgray!18] (-2.1,-1.5) rectangle (2.1,0);
  \draw[figline] (-2.1,0) -- (2.1,0);
  \node[figlblaux] at (0,-1.15) {conductor: $\vec E=0$};
  % densidad superficial sobre el borde
  \foreach \xx in {-0.72,-0.43,-0.14,0.14,0.43,0.72}{\node[figpos] at (\xx,0) {};}
  \node[figlbl, figred] at (1.55,-0.3) {$\sigma$};
  \draw[figvecaux] (1.42,-0.24) -- (0.9,-0.04);
  % cilindrito chato
  \draw[figred, line width=1pt] (-0.9,0.24) -- (0.9,0.24) -- (0.9,-0.24) -- (-0.9,-0.24) -- cycle;
  \node[figlbl, figred] at (1.2,0.35) {$S$};
  % tapa exterior: la unica que aporta flujo
  \draw[figvecres] (-0.45,0.24) -- (-0.45,1.1);
  \draw[figvecres] (0.45,0.24) -- (0.45,1.1);
  \node[figlbl] at (0,1.3) {$E_n$};
  % cota de la tapa
  \draw[figaux] (-0.9,-0.42) -- (-0.9,-0.86);
  \draw[figaux] (0.9,-0.42) -- (0.9,-0.86);
  \draw[figvecaux] (-0.9,-0.68) -- (0.9,-0.68);
  \draw[figvecaux] (0.9,-0.68) -- (-0.9,-0.68);
  \node[figlbl, fill=figgray!18, inner sep=1.5pt] at (0,-0.68) {$\Delta S$};
  \node[figlbl] at (0,-2.25) {(b) $E_n\,\Delta S=\sigma\,\Delta S/\varepsilon_0$};
  \node[figlbl] at (0,-2.85) {$\Rightarrow\ E_n=\sigma/\varepsilon_0$};
\end{scope}

\end{tikzpicture}

  \caption*{\small\itshape Los dos pasos, en orden: la circulación mata la
  componente tangencial sin usar Gauss (a), y recién entonces el cilindrito
  gaussiano da la normal (b).}
\end{figure}

\begin{itemize}[nosep]
  \item La tapa interior no aporta: $\vec E = 0$.
  \item La cara lateral no aporta: es de altura despreciable y, además, el campo
        exterior es normal.
  \item La tapa exterior aporta $E_n\,\Delta S$.
\end{itemize}

La carga encerrada está concentrada en la superficie, así que vale
$\sigma\,\Delta S$. Gauss da $E_n\,\Delta S = \sigma\,\Delta S/\varepsilon_0$, o
sea

\[
\boxed{\;E_n = \frac{\sigma}{\varepsilon_0}\;}
\]

\begin{notebox}[Dos precisiones sobre este resultado]
\begin{enumerate}[nosep, leftmargin=*]
  \item \textbf{No hace falta que $\sigma$ sea uniforme.} En una superficie de
        forma cualquiera la densidad se distribuye de manera inhomogénea, pero la
        relación vale \textbf{punto a punto}: en cada punto el campo es normal y
        su valor es el $\sigma$ \emph{de ese punto} sobre $\varepsilon_0$.
  \item \textbf{El signo importa.} La fórmula tal como está da la componente
        normal saliente con su signo: si $\sigma>0$ el campo apunta hacia afuera,
        y si $\sigma<0$ apunta hacia adentro. Si se escribe la igualdad entre
        módulos hay que poner $|\vec E| = |\sigma|/\varepsilon_0$.
\end{enumerate}
\end{notebox}

\section{La idea del desarrollo multipolar}

Con esto termina el repaso y empieza el tema nuevo. La motivación es una pregunta
simple: si se tiene una \textbf{distribución localizada} de carga ---discreta o
continua, pero confinada a una región--- y se la mira desde muy lejos, ¿cómo se
comporta?

\begin{itemize}[nosep]
  \item \textbf{Casi siempre, como una carga puntual}, con la carga neta total.
  \item \textbf{Si la carga neta es nula}, ese efecto desaparece y lo que queda,
        visto de lejos, es un objeto un poco más complicado: un \textbf{dipolo
        eléctrico}.
  \item \textbf{Si además el momento dipolar es nulo}, lo que queda es un objeto
        todavía más complicado: un \textbf{cuadrupolo}.
  \item Y así sucesivamente.
\end{itemize}

\begin{keybox}[La idea de fondo]
Toda distribución localizada de carga, vista de lejos, \textbf{se simplifica}: se
comporta como una jerarquía de objetos caracterizables, y uno puede llevar la
descripción al nivel de precisión que quiera agregando términos.
\end{keybox}

El orden de la clase es deliberado: primero el dipolo (que ya se vio en Física
III, aunque no tan matemáticamente), como calentamiento; después la
generalización a una distribución cualquiera.

\section{El dipolo eléctrico}

\subsection{Campo exacto de dos cargas opuestas}

Un \textbf{dipolo}, en esta primera definición, son dos cargas de igual módulo y
signo opuesto. (El docente advierte que ésta no es la única definición y que más
adelante vendrá una más general.)

Geometría: con el origen fijado, la carga $-Q$ está en $\vec r\,'$, y la carga
$+Q$ en $\vec r\,' + \vec L$; el vector $\vec L$ va de la negativa a la positiva.
El punto de observación $P$ está en $\vec r$.

El campo es pura superposición de dos Coulomb, sin ninguna aproximación:

\[
\vec E(\vec r) = \frac{Q}{4\pi\varepsilon_0}\left[
  \frac{\vec r - \vec r\,' - \vec L}{|\vec r - \vec r\,' - \vec L|^3}
  - \frac{\vec r - \vec r\,'}{|\vec r - \vec r\,'|^3}\right]
\]

Hasta acá, en palabras del docente, «es una trivialidad, es simplemente escribir
Coulomb». Lo que interesa ---y no es trivial--- es el comportamiento a
\textbf{gran distancia}.

\begin{figure}[H]
  \centering
  % Clase 4 — geometria del dipolo: -Q en r', +Q en r'+L, punto de observacion P en r.
% Es el dibujo que el docente conserva en el pizarron durante todo el desarrollo.
% Los vectores r' y r se dibujan con angulo bien separado a proposito: si quedan
% casi colineales, r', L, r y r-r' se encinan y la figura no se lee.
\begin{tikzpicture}[scale=1.15]

  % origen
  \node[figneutra] at (0,0) {};
  \node[figlbl] at (-0.32,-0.2) {$O$};

  % cargas del dipolo, abajo a la derecha del origen
  \node[figneg] (mq) at (1.85,-0.95) {};
  \node[figpos] (pq) at (2.72,-0.42) {};
  \node[figlbl, figblue] at (1.72,-1.32) {$-Q$};
  \node[figlbl, figred]  at (3.02,-0.72) {$+Q$};

  % vectores de posicion de cada carga
  \draw[figvec] (0,0) -- (mq);
  \node[figlbl, fill=white, inner sep=1.5pt] at (0.85,-0.62) {$\vec r\,'$};
  \draw[figvecaux] (0,0) -- (pq);
  \node[figlblaux, fill=white, inner sep=1.5pt] at (1.12,0.05) {$\vec r\,'+\vec L$};

  % vector L, de la negativa a la positiva
  \draw[figvecres] (mq) -- (pq);
  \node[figlbl, figred, fill=white, inner sep=1.5pt] at (2.34,-0.95) {$\vec L$};

  % punto de observacion, bien separado angularmente
  \node[figneutra] (P) at (4.35,2.75) {};
  \node[figlbl] at (4.6,2.95) {$P$};
  \draw[figvec] (0,0) -- (P);
  \node[figlbl, fill=white, inner sep=1.5pt] at (1.75,1.35) {$\vec r$};

  % separacion r - r'
  \draw[figline, dashed] (mq) -- (P);
  \node[figlbl, fill=white, inner sep=1.5pt] at (3.62,1.15) {$\vec r-\vec r\,'$};

  \node[figlbl, align=center] at (2.2,-2.35)
    {Desarrollo de Taylor para $L\ll|\vec r-\vec r\,'|$;\\[0.2em]
     momento dipolar $\vec p=Q\vec L$};

\end{tikzpicture}

  \caption*{\small\itshape La geometría que el docente conserva en el pizarrón
  durante todo el desarrollo: el origen es libre, y lo que se compara es $L$
  contra $|\vec r-\vec r\,'|$.}
\end{figure}

\subsection{La fórmula maestra del desarrollo}

«Gran distancia» significa $|\vec r - \vec r\,'| \gg L$. Antes de desarrollar hay
dos observaciones metodológicas que el docente subraya:

\begin{enumerate}[leftmargin=*]
  \item \textbf{Hacer $|\vec r - \vec r\,'|$ grande es lo mismo que hacer
        $L\to 0$}: lo que importa es la comparación entre las dos, no el valor de
        ninguna. Y el límite $L\to 0$ es trivial (da cero); lo interesante es la
        \textbf{corrección} a ese límite, es decir, el desarrollo de Taylor.
  \item \textbf{Los desarrollos se hacen en cantidades adimensionadas.} No tiene
        sentido decir que una distancia es grande o chica: es chica
        \emph{comparada con otra}. El parámetro del desarrollo es
        $L/|\vec r - \vec r\,'|$.
\end{enumerate}

Conviene entonces resolver \textbf{una sola vez} la expansión genérica que
después se usará con distintos exponentes. Sea $n$ entero. El truco es escribir
la potencia como un módulo al cuadrado, para poder usar que el módulo al cuadrado
de un vector es su producto escalar consigo mismo:

\[
\frac{1}{|\vec r - \vec r\,' - \vec L|^{n}}
= \left(|\vec r - \vec r\,' - \vec L|^2\right)^{-n/2}
= \left(|\vec r - \vec r\,'|^2 + L^2 - 2\,\vec L\cdot(\vec r - \vec r\,')\right)^{-n/2}
\]

Se factoriza el orden cero para dejar el corchete adimensionado:

\[
= \frac{1}{|\vec r - \vec r\,'|^{n}}
\left[1 + \frac{L^2}{|\vec r - \vec r\,'|^2}
        - \frac{2\,\vec L\cdot(\vec r - \vec r\,')}{|\vec r - \vec r\,'|^2}\right]^{-n/2}
\]

Dentro del corchete el primer término correctivo es de orden
$(L/|\vec r-\vec r\,'|)^2$ y el segundo de orden $L/|\vec r-\vec r\,'|$: a primer
orden sólo sobrevive el segundo. Falta el desarrollo de $f(x) = (1+x)^\alpha$, que
se hace «a pedal»: $f(0)=1$, $f'(x) = \alpha(1+x)^{\alpha-1}$, $f'(0)=\alpha$, de
donde

\[
(1+x)^{\alpha} = 1 + \alpha x + \mathcal{O}(x^2)
\]

Con $\alpha = -n/2$ y
$x = -2\vec L\cdot(\vec r-\vec r\,')/|\vec r-\vec r\,'|^2$, los dos signos menos y
el 2 se cancelan contra el $1/2$:

\[
\boxed{\;\frac{1}{|\vec r - \vec r\,' - \vec L|^{n}}
= \frac{1}{|\vec r - \vec r\,'|^{n}}
  \left[1 + n\,\frac{\vec L\cdot(\vec r - \vec r\,')}{|\vec r - \vec r\,'|^{2}}\right]
  + \mathcal{O}(L^2)\;}
\]

Esta fórmula se usa dos veces en la clase: con $n=3$ para el campo y con $n=1$
para el potencial.

\subsection{Campo del dipolo a gran distancia}

Aplicando la fórmula maestra con $n = 3$ al primer término del campo, y notando
que el factor $(\vec r - \vec r\,' - \vec L)$ del numerador ya contiene un
$\vec L$ ---de modo que multiplicarlo por la corrección lineal daría
$\mathcal{O}(L^2)$---:

\[
\frac{\vec r - \vec r\,' - \vec L}{|\vec r - \vec r\,' - \vec L|^3}
= \frac{\vec r - \vec r\,'}{|\vec r - \vec r\,'|^3}
- \frac{\vec L}{|\vec r - \vec r\,'|^3}
+ \frac{3\,[\vec L\cdot(\vec r - \vec r\,')]\,(\vec r - \vec r\,')}{|\vec r - \vec r\,'|^5}
+ \mathcal{O}(L^2)
\]

Al restar el campo de la carga negativa, \textbf{el término de orden cero se
cancela} --- y tenía que cancelarse:

\begin{notebox}[Por qué el orden cero tiene que irse]
Si el término de orden cero no se fuera, el sistema visto de lejos sería
equivalente a una carga puntual. Pero su carga neta es cero, así que el efecto de
carga puntual debe anularse. «Si no se fuera, cometimos un error.»
\end{notebox}

Definiendo el \textbf{momento dipolar eléctrico}

\[
\boxed{\;\vec p = Q\,\vec L\;}
\]

queda, para $L \ll |\vec r - \vec r\,'|$,

\[
\boxed{\;\vec E_{\text{dip}}(\vec r) \simeq \frac{1}{4\pi\varepsilon_0}\left[
  \frac{3\,[\vec p\cdot(\vec r - \vec r\,')]\,(\vec r - \vec r\,')}{|\vec r - \vec r\,'|^{5}}
  - \frac{\vec p}{|\vec r - \vec r\,'|^{3}}\right] + \mathcal{O}(L^2)\;}
\]

\begin{warnbox}[La letra $p$]
Es la misma que la de cantidad de movimiento y \textbf{no tiene nada que ver} con
ella.
\end{warnbox}

Las correcciones descartadas, de orden $L^2$, se llaman \textbf{correcciones
cuadrupolares}. No se calculan acá: el docente aclara que se obtienen llevando el
mismo desarrollo de Taylor un orden más lejos, y que eso es justamente lo que
hará después para una distribución cualquiera.

\subsection{Potencial del dipolo}

El mismo cálculo para el potencial es más corto, y se hace directamente en vez de
integrar el campo. Superponiendo los dos potenciales de Coulomb:

\[
\phi(\vec r) = \frac{Q}{4\pi\varepsilon_0}\left[
  \frac{1}{|\vec r - \vec r\,' - \vec L|} - \frac{1}{|\vec r - \vec r\,'|}\right]
\]

Ahora la fórmula maestra se usa con $n = 1$. De nuevo el orden cero se cancela
contra el segundo término y queda

\[
\boxed{\;\phi_{\text{dip}}(\vec r) \simeq
  \frac{1}{4\pi\varepsilon_0}\,\frac{\vec p\cdot(\vec r - \vec r\,')}{|\vec r - \vec r\,'|^{3}}\;}
\]

una expresión notoriamente más simple que la del campo.

\begin{notebox}[Ejercicio dejado en clase]
Verificar que $\vec E_{\text{dip}} = -\grad\phi_{\text{dip}}$ con las dos
expresiones obtenidas. Dan exactamente lo mismo porque en ambos cálculos se
descartaron los mismos términos de orden $L^2$ y ambos son lineales en $\vec L$.
\end{notebox}

\section{Dipolo en un campo externo: energía y fuerza}

El problema se da vuelta: ya no interesa el campo \emph{producido} por el dipolo,
sino la fuerza que \emph{experimenta} cuando se lo coloca en un campo externo
$\vec E_{\text{ext}}$ generado por otras cargas.

\textbf{Primera observación, antes de toda cuenta: si el campo externo es
uniforme, la fuerza neta es cero.} La fuerza sobre $+Q$ y la fuerza sobre $-Q$
tienen igual módulo y dirección y sentidos opuestos, de modo que se cancelan. (Hay
un torque neto, que el docente menciona pero no calcula.)

\begin{figure}[H]
  \centering
  % Clase 4 — dipolo rigido ("barrita") en un campo externo.
% (a) campo uniforme: las fuerzas se cancelan, fuerza neta nula (queda torque).
% (b) campo inhomogeneo: los modulos difieren y aparece fuerza neta.
% Los rotulos de carga van ARRIBA/ABAJO de la carga y los de fuerza PEGADOS a su
% flecha: puestos a la misma altura se leen como un solo bloque (pasa en el render,
% no en el codigo).
\begin{tikzpicture}[scale=1]

% ---------------- panel (a) ----------------
\begin{scope}
  \foreach \yy in {-1.65,-0.95,0.95,1.65}{
    \draw[figvecaux] (-2.4,\yy) -- (2.4,\yy);
  }
  \node[figlblaux] at (0,2.1) {$\vec E_{\text{ext}}$ uniforme};
  \draw[figline, line width=1.4pt] (-0.7,-0.42) -- (0.7,0.42);
  \node[figneg] (a) at (-0.7,-0.42) {};
  \node[figpos] (b) at (0.7,0.42) {};
  \node[figlbl, figred,  fill=white, inner sep=1.5pt] at (0.62,0.8)  {$+Q$};
  \node[figlbl, figblue, fill=white, inner sep=1.5pt] at (-0.62,-0.8) {$-Q$};
  \node[figlbl, fill=white, inner sep=1.5pt] at (-0.08,0.16) {$\vec L$};
  % fuerzas: igual modulo, sentidos opuestos
  \draw[figvecres] (b) -- ++(1.2,0);
  \draw[figvecres] (a) -- ++(-1.2,0);
  \node[figlbl, fill=white, inner sep=1.5pt] at (1.5,0.14)  {$Q\vec E$};
  \node[figlbl, fill=white, inner sep=1.5pt] at (-1.5,-0.14) {$-Q\vec E$};
  \node[figlbl] at (0,-2.6) {(a) fuerza neta $=0$};
\end{scope}

% ---------------- panel (b) ----------------
\begin{scope}[xshift=7.4cm]
  \draw[figvecaux] (-2.4,1.65) -- (2.4,0.9);
  \draw[figvecaux] (-2.4,0.95) -- (2.4,0.52);
  \draw[figvecaux] (-2.4,-0.95) -- (2.4,-0.52);
  \draw[figvecaux] (-2.4,-1.65) -- (2.4,-0.9);
  \node[figlblaux] at (0,2.1) {$\vec E_{\text{ext}}$ inhomog\'eneo};
  \draw[figline, line width=1.4pt] (-0.7,-0.42) -- (0.7,0.42);
  \node[figneg] (c) at (-0.7,-0.42) {};
  \node[figpos] (d) at (0.7,0.42) {};
  \node[figlbl, figred,  fill=white, inner sep=1.5pt] at (0.62,0.8)  {$+Q$};
  \node[figlbl, figblue, fill=white, inner sep=1.5pt] at (-0.62,-0.8) {$-Q$};
  % fuerzas de distinto modulo: no se cancelan
  \draw[figvecres] (d) -- ++(1.6,0);
  \draw[figvecres] (c) -- ++(-0.6,0);
  \node[figlbl, fill=white, inner sep=1.5pt] at (1.7,0.14)  {$Q\vec E$};
  \node[figlbl, fill=white, inner sep=1.5pt] at (-1.28,-0.14) {$-Q\vec E$};
  \node[figlbl] at (0,-2.6) {(b) fuerza neta $\neq0$};
\end{scope}

\node[figlbl, align=center] at (3.7,-3.55)
  {$U=-\vec p\cdot\vec E_{\text{ext}}$,\quad $\vec F=-\nabla U$.\ \ La aproximaci\'on pide\\[0.2em]
   $L\ll H$, con $H$ la escala en que var\'ia $\vec E_{\text{ext}}$};

\end{tikzpicture}

  \caption*{\small\itshape La misma barrita en los dos regímenes: lo único que
  cambia entre los paneles es la homogeneidad del campo externo, y con ella el
  balance de las dos fuerzas.}
\end{figure}

\begin{keybox}
Para que haya \textbf{fuerza} sobre un dipolo, el campo externo tiene que ser
\textbf{inhomogéneo}.
\end{keybox}

Conviene trabajar con el potencial externo, que es más simple. Modelando el
dipolo como una barrita rígida de longitud $L$, con $-Q$ en $\vec r$ y $+Q$ en
$\vec r + \vec L$, la energía es la suma de carga por potencial de cada una:

\[
U = Q\,\phi_{\text{ext}}(\vec r + \vec L) - Q\,\phi_{\text{ext}}(\vec r)
\]

\textbf{¿Chico respecto de qué?} Acá el docente se detiene en una confusión
frecuente: $L$ \textbf{no} puede ser chico «respecto de $\vec r$», porque
$\vec r$ depende de dónde se ponga el origen, que es una convención (se podría
elegir $\vec r = 0$). La escala de comparación correcta es la \textbf{distancia
típica $H$ en la que el campo externo varía apreciablemente}: la aproximación es
$L \ll H$. Es coherente con lo anterior --- si el campo es uniforme,
$H\to\infty$ y no hay efecto.

El desarrollo de Taylor es ahora \textbf{de una función de varias variables}, y se
hace por coordenadas:

\[
\phi_{\text{ext}}(x + L_x,\, y + L_y,\, z + L_z)
= \phi_{\text{ext}}(\vec r)
+ \pdv{\phi_{\text{ext}}}{x}L_x
+ \pdv{\phi_{\text{ext}}}{y}L_y
+ \pdv{\phi_{\text{ext}}}{z}L_z
+ \mathcal{O}(L^2)
\]

Los tres términos lineales son un producto escalar del gradiente con $\vec L$:

\[
\phi_{\text{ext}}(\vec r + \vec L)
= \phi_{\text{ext}}(\vec r) + \vec L\cdot\grad\phi_{\text{ext}}(\vec r)
+ \mathcal{O}(L^2)
\]

Sustituyendo, \textbf{otra vez el orden cero se cancela} ---por la misma razón de
siempre: si no hubiera inhomogeneidad no habría efecto---, y con
$\vec p = Q\vec L$:

\[
U = \vec p\cdot\grad\phi_{\text{ext}}(\vec r) + \mathcal{O}(L^2)
\]

Finalmente, como $\grad\phi = -\vec E$:

\[
\boxed{\;U = -\,\vec p\cdot\vec E_{\text{ext}}(\vec r) + \mathcal{O}(L^2)\;}
\]

\begin{warnbox}[El signo]
$\grad\phi = -\vec E$, no $+\vec E$. Es un signo fácil de olvidar y que no se debe
olvidar: se lo lleva puesto toda la fórmula de la energía.
\end{warnbox}

La fuerza se obtiene como siempre a partir de la energía potencial,
$\vec F = -\grad U$. Se ve inmediatamente la consistencia con la observación
inicial: si $\vec E_{\text{ext}}$ es uniforme, $U$ no depende de $\vec r$ y su
gradiente es nulo, de modo que la fuerza es cero.

El término $\mathcal{O}(L^2)$ descartado es el que corresponde a \textbf{dipolos
grandes} ---grandes frente a la escala de variación del campo---. Que despreciarlo
sea buena aproximación depende del nivel de precisión que se busque; la fórmula
general vale siempre, la simplificación es la que tiene condiciones.

\section{Desarrollo multipolar de una distribución continua}

El último tramo generaliza todo lo anterior de dos cargas discretas a una
\textbf{distribución continua cualquiera}. (El cálculo para una distribución
discreta es casi idéntico.)

\textbf{Planteo.} Una densidad volumétrica $\rho(\vec r\,')$ localizada en un
volumen $B$: \emph{localizada} significa que fuera de un volumen finito es
despreciable. Se toma el origen \textbf{dentro} de $B$ ---en un punto cualquiera,
no importa cuál--- y se llama $a$ al \textbf{tamaño típico} de $B$ (por ejemplo,
el diámetro de una esfera que lo contenga). El punto de observación $P$ está en
$\vec r$, con $r \gg a$.

El potencial exacto lo da Coulomb, sin aproximación alguna:

\[
\phi(\vec r) = \frac{1}{4\pi\varepsilon_0}\int_B \frac{\rho(\vec r\,')}{|\vec r - \vec r\,'|}\,dV'
\]

donde la prima en $dV' = dx'\,dy'\,dz'$ insiste en que la integración es sobre
$\vec r\,'$.

\begin{figure}[H]
  \centering
  % Clase 4 — planteo del desarrollo multipolar: distribucion localizada rho(r') en
% un volumen B de tamano tipico a, origen dentro de B, observacion en P con r >> a.
% Ojo: r y r-r' son casi colineales por construccion (r >> a), asi que sus rotulos
% se separan perpendicularmente, uno debajo de la flecha y otro encima de la punteada.
\begin{tikzpicture}[scale=1.15]

  % volumen B
  \fill[figblue!12] plot[smooth cycle, tension=0.85]
    coordinates {(-1.0,0.5) (-0.1,1.05) (0.95,0.6) (1.05,-0.4) (0.1,-0.95) (-0.95,-0.45)};
  \draw[figline] plot[smooth cycle, tension=0.85]
    coordinates {(-1.0,0.5) (-0.1,1.05) (0.95,0.6) (1.05,-0.4) (0.1,-0.95) (-0.95,-0.45)};
  \node[figlbl] at (-1.5,1.15) {$B$};
  \node[figlbl] at (0.42,-1.85) {$\rho(\vec r\,')$};
  \draw[figvecaux] (0.42,-1.62) -- (0.32,-0.62);

  % esfera de tamano tipico a: la cota va a la izquierda, lejos de r' y de rho
  \draw[figaux] (0.02,0.02) circle (1.22);
  \draw[figvecaux] (0.02,0.02) -- (-1.2,0.02);
  \node[figlblaux, fill=figblue!12, inner sep=1.5pt] at (-0.62,0.02) {$a$};

  % origen dentro de B
  \node[figneutra] at (0,0) {};
  \node[figlbl, fill=figblue!12, inner sep=1.5pt] at (-0.02,-0.34) {$O$};

  % punto fuente dentro de la distribucion
  \node[figneutra] (s) at (0.6,0.52) {};
  \draw[figvec] (0,0) -- (s);
  \node[figlbl, fill=figblue!12, inner sep=1.5pt] at (0.16,0.62) {$\vec r\,'$};

  % punto de observacion, lejos
  \node[figneutra] (P) at (5.6,2.3) {};
  \node[figlbl] at (5.85,2.5) {$P$};
  \draw[figvec] (0,0) -- (P);
  \node[figlbl, fill=white, inner sep=1.5pt] at (2.3,0.62) {$\vec r$};
  \draw[figaux] (2.3,0.8) -- (2.3,0.93);

  % separacion, casi colineal con r: su rotulo va del otro lado
  \draw[figline, dashed] (s) -- (P);
  \node[figlblaux, fill=white, inner sep=1.5pt] at (3.6,2.05) {$|\vec r-\vec r\,'|$};
  \draw[figaux] (3.6,1.88) -- (3.6,1.66);

  \node[figlbl, align=center] at (2.6,-2.75)
    {Distribuci\'on localizada: $\rho$ despreciable fuera de $B$.\\[0.2em]
     El desarrollo es en $r'/r$, acotado por $a/r\ll1$};

\end{tikzpicture}

  \caption*{\small\itshape El origen se elige \emph{dentro} de $B$ y su ubicación
  exacta no importa: lo que gobierna el desarrollo es el cociente entre el tamaño
  típico $a$ y la distancia de observación $r$.}
\end{figure}

\textbf{Toda la dependencia en $\vec r$ está en el denominador}, así que el
desarrollo se hace ahí. Desarrollar en $r$ grande es lo mismo que desarrollar en
$r' \ll r$, porque $r'$ está acotado por $a$.

Misma estrategia que en la sección anterior: escribir el módulo como producto
escalar y factorizar el orden cero.

\[
\frac{1}{|\vec r - \vec r\,'|}
= \left(r^2 - 2\,\vec r\cdot\vec r\,' + r'^2\right)^{-1/2}
= \frac{1}{r}\left[1 - \frac{2\,\vec r\cdot\vec r\,'}{r^2} + \frac{r'^2}{r^2}\right]^{-1/2}
\]

Pero ahora \textbf{no alcanza el primer orden}: hay que ir al segundo, porque el
objetivo es llegar hasta el término cuadrupolar. Se retoma $f(x)=(1+x)^\alpha$ y
se calcula una derivada más:

\[
f''(x) = \alpha(\alpha-1)(1+x)^{\alpha-2}
\quad\Longrightarrow\quad
f''(0) = \alpha(\alpha-1)
\]

\[
(1+x)^{\alpha} = 1 + \alpha x + \tfrac{1}{2}\,\alpha(\alpha-1)\,x^2 + \mathcal{O}(x^3)
\]

\begin{notebox}[El factor $1/n!$]
Es el que da el $1/2$; el término siguiente sería
$\tfrac{1}{6}\alpha(\alpha-1)(\alpha-2)x^3$. Es fácil ir a cualquier orden.
\end{notebox}

Con $\alpha = -1/2$ y $x = -2\,\vec r\cdot\vec r\,'/r^2 + r'^2/r^2$:

\begin{itemize}[leftmargin=*]
  \item del término $\alpha x$ salen
        $+\dfrac{\vec r\cdot\vec r\,'}{r^2}$ y
        $-\dfrac{1}{2}\dfrac{r'^2}{r^2}$;
  \item de $\tfrac12\alpha(\alpha-1)x^2 = \tfrac{3}{8}x^2$ hay que
        \textbf{quedarse sólo con el cuadrado del término cruzado}: al elevar $x$
        al cuadrado aparecen un término de orden $r'^2$, otro de orden $r'^4$ y
        un cruzado de orden $r'^3$; los dos últimos ya son de orden superior. El
        que sobrevive es
        $\tfrac{3}{8}\cdot\dfrac{4(\vec r\cdot\vec r\,')^2}{r^4}
         = \dfrac{3}{2}\dfrac{(\vec r\cdot\vec r\,')^2}{r^4}$.
\end{itemize}

Resultado del desarrollo, que es donde queda la clase:

\[
\boxed{\;\frac{1}{|\vec r - \vec r\,'|}
= \frac{1}{r}\left[1 + \frac{\vec r\cdot\vec r\,'}{r^{2}}
                     - \frac{1}{2}\frac{r'^{2}}{r^{2}}
                     + \frac{3}{2}\frac{(\vec r\cdot\vec r\,')^{2}}{r^{4}}\right]
  + \mathcal{O}\!\left(\left(\tfrac{r'}{r}\right)^{3}\right)\;}
\]

\begin{notebox}[Sobre el coeficiente del último término]
El $3/2$ se armó en el pizarrón en dos pasos y con una corrección sobre la
marcha: al factor $3/8$ hay que multiplicarlo por el $4$ que sale de elevar al
cuadrado el $-2$ del término cruzado. El propio docente detectó en clase que se le
había «morfado» un $1/2$ intermedio y lo corrigió.
\end{notebox}

La clase termina acá, deliberadamente en suspenso: falta meter este desarrollo
dentro de la integral y ver qué momentos de la distribución multiplican a cada
potencia de $1/r$. Eso es el desarrollo multipolar propiamente dicho.

\bigskip
\noindent\emph{Continúa en la Clase 5, donde el desarrollo se sustituye en la
integral y aparecen la carga total, el momento dipolar y el tensor cuadrupolar,
para luego pasar a las ecuaciones de Poisson y Laplace.}

\end{document}
